% Part: computability
% Chapter: tm-computations
% Section: church-turing-thesis

\documentclass[../../../include/open-logic-section]{subfiles}

\begin{document}

\olfileid{tur}{mac}{ctt}
\olsection{चर्च--ट्यूरिंग प्रबंध}

प्रभावी प्रक्रियेच्या संकल्पनेला अचूक पर्याय म्हणून ट्यूरिंग
यंत्रे मांडली जातात. ट्यूरिंग यांच्या मते, प्रभावी प्रक्रिया आणि
ट्यूरिंग यंत्र या दोन्ही संकल्पना समजलेल्या कोणालाही प्रभावी
प्रक्रियेने करता येणारी कोणतीही गोष्ट ट्यूरिंग यंत्रानेही करता
येईल असे अंतर्ज्ञान होईल. प्रभावी प्रक्रियेच्या संकल्पनेचे
इतर सर्व प्रस्तावित अचूक पर्याय व्याप्तीच्या दृष्टीने ट्यूरिंग
यंत्राच्या संकल्पनेशी सममूल्य ठरतात, ही गोष्ट या दाव्याला
आधार देते---म्हणजे ते नेमका तोच फलनांचा संच संगणित करू
शकतात. या दाव्याला \emph{चर्च--ट्यूरिंग प्रबंध} म्हणतात.

\begin{defn}[चर्च--ट्यूरिंग प्रबंध]
\emph{चर्च--ट्यूरिंग प्रबंध} असे सांगतो की प्रभावी प्रक्रियेने
संगणित करता येणारी कोणतीही गोष्ट ट्यूरिंग-संगणनक्षम असते.
\end{defn}

चर्च--ट्यूरिंग प्रबंधाचा दोन प्रकारे आधार घेतला जातो. पहिला
वापर काहीसा आळशीपणाला सबब देणारा आहे. समजा, एखादी गोष्ट
संगणित करण्याच्या प्रभावी प्रक्रियेचे वर्णन आपल्याकडे
``छद्म-संकेतलेखनात'' आहे. अशा वेळी प्रत्यक्षात ट्यूरिंग यंत्र
बनवले नसले तरी हेच फलन एखादे ट्यूरिंग यंत्र संगणित करते,
असा दावा चर्च--ट्यूरिंग प्रबंधाच्या आधारे समर्थित करता येतो.

चर्च--ट्यूरिंग प्रबंधाचा दुसरा वापर तत्त्वज्ञानाच्या दृष्टीने
अधिक रोचक आहे. काही फलने ट्यूरिंग यंत्रांनी संगणित करता
येत नाहीत, हे दाखवता येते. त्यावरून चर्च--ट्यूरिंग प्रबंध
वापरून कोणत्याही प्रक्रियेने ती प्रभावी रीतीने संगणित करता
येत नाहीत, असा निष्कर्ष काढता येतो. कारण तशी एखादी प्रक्रिया
असती तर चर्च--ट्यूरिंग प्रबंधानुसार त्यासाठी ट्यूरिंग यंत्रही
असले असते. म्हणून एखादे फलन संगणित करणारे ट्यूरिंग यंत्र
असू शकत नाही हे सिद्ध केले, तर त्यासाठी प्रभावी प्रक्रियाही
असू शकत नाही. विशेषतः, तथाकथित थांबण्याची समस्या केवळ
ट्यूरिंग यंत्रांनीच सोडवता येत नाही असे नव्हे; ती मुळीच
प्रभावी रीतीने सोडवता येत नाही, असा दावा करण्यासाठी
चर्च--ट्यूरिंग प्रबंध वापरला जातो.

\end{document}
